%% ════════════════════════════════════════════════════════════════════════════
\section{Data and Preprocessing}
\label{sec:data}

\subsection{Data Sources and Provenance}
We integrate two official, open data repositories published under the statutory framework of India's CPCB:
\begin{enumerate}
  \item \textbf{CPCB Daily Air Quality Bulletins (2015--2023):} Official 4:00~PM bulletin reports published daily across all monitored municipal centres, archived on Zenodo (DOI: \href{https://doi.org/10.5281/zenodo.11146316}{10.5281/zenodo.11146316}; 299,974 records). Each bulletin reports the official municipal ``Index Value'', calculated according to CPCB guidelines based on active CAAQMS stations. This official index forms our primary target variable $\yt$.
  \item \textbf{CPCB Station Concentrations (Kaggle CAAQMS Archive):} Daily ambient concentration measurements for the seven criteria pollutants ($PM_{2.5}, PM_{10}, NO_2, SO_2, CO, O_3, NH_3$) spanning 2015--2020. Where station-aggregated AQI values diverged from official bulletins due to unbounded sub-index equations exceeding the theoretical maximum of 500, the official bulletin figure was enforced as the authoritative ground truth.
\end{enumerate}

\subsection{City Qualification and Regional Scope}
To evaluate geographical generalization, six municipal centres were qualified based on temporal continuity, data density, and industrial relevance. Table~\ref{tab:city_qualification} details the qualified cities, monitoring intervals, record counts, and assigned experimental roles.

\input{../results/tables/table1_qual.tex}

Ahmedabad and Chennai satisfy the longest historical continuity ($n = 3{,}167$ calendar days) and possess concurrent criteria pollutant measurements, qualifying for all three comparative setups (A, B, and C). Ahmedabad represents inland semi-arid western India, whereas Chennai serves as a maritime southern comparative anchor. Gandhinagar, Ankleshwar, and Vapi qualify from 2019 onward based on official CPCB daily bulletins. Surat, where CAAQMS stations first appear in the CPCB bulletin archive in late 2021 ($n = 423$ bulletin days; 327 valid observations with non-missing AQI), serves as a thin-data sensitivity evaluation. (The precise commissioning date of Surat's CAAQMS network is not formally documented in the publicly available CPCB technical reports; we infer it from the first occurrence of valid Surat entries in the Zenodo bulletin archive.)

\subsection{Data Cleaning and Leakage Prevention}
The raw municipal series were standardised and processed through rigorous quality control:
\begin{itemize}
  \item \textbf{Physical Boundary Filtering:} Negative values and physically implausible sensor readings outside CPCB technical ranges were flagged and set to missing. In Ahmedabad, ten spurious sensor readings for CO exceeding $100\,\text{mg/m}^3$ were removed.
  \item \textbf{Continuous Calendar Reindexing:} All city time series were reindexed to a continuous daily calendar. Missing observations in the target variable $\yt$ were recorded as unobserved; they were strictly excluded from loss computation and were never synthetically interpolated.
  \item \textbf{Strict Pre-Split Feature Imputation:} To guarantee zero lookahead bias, missing input feature values were imputed using training-partition medians only. Imputation statistics were calculated strictly on the training partition and subsequently applied to test observations.
\end{itemize}
All validated datasets are archived in binary Parquet format with cryptographic SHA-256 integrity checksums registered in the project manifest.

